\section{Endpoint margins and uniform coefficient estimates}\label{sec:repairs}

Two estimates in the proof of Tao's Proposition~5.2 control the location
of the first-passage time and the effect of fine-scale mixing on the
first-entry distribution.  Retain his parameters
\begin{equation}\label{eq:local-parameters}
 L=\log x,\quad \alpha=1.001,\quad \beta=\log(4/3),\quad
 n_0=\left\lfloor\frac{L}{10\log2}\right\rfloor,\quad
 m_0=\left\lfloor\frac{L}{100000}\right\rfloor.
\end{equation}
Here $x$ is sufficiently large; in particular $m_0\geq1$.

\subsection{A buffer in time is multiplicative in the initial value}

In the proof of Proposition~5.2, Tao obtains an event $G_x$ with
$\Prob(G_x^c)=O(L^{-10})$ on which, for
$\mathbf N_y\sim\mu_y$ and $y\in\{x^\alpha,x^{\alpha^2}\}$,
\begin{equation}\label{eq:orbit-log}
 \left|\log\frac{\Syr^j(\mathbf N_y)}{\mathbf N_y}
          +j\beta\right|\leq C_0L^{3/5}
 \quad(0\leq j\leq n_0).
\end{equation}
The desired time interval is
\[
 I_y=\left[\frac{\log(y/x)}\beta+L^{4/5},\quad
             \frac{\log(y^\alpha/x)}\beta-L^{4/5}\right].
\]
The input restriction printed at this point in v7 adds and subtracts
$2L^{4/5}$ at the endpoints $y,y^\alpha$.
An additive displacement of that size produces logarithmic displacement
$O(L^{4/5}/y)$ at the lower endpoint, not the required time margin.
A multiplicative restriction supplies the required margin.

\begin{proposition}[Multiplicative endpoint estimate]\label{prop:buffer}
Set
\[
 J_y=[y\exp(L^{4/5}),\ y^\alpha\exp(-L^{4/5})].
\]
For all sufficiently large $x$ this is nonempty,
\[
 G_x\cap\{\mathbf N_y\in J_y\}\subseteq\{t_x(\mathbf N_y)\in I_y\},
 \qquad
 \Prob(t_x(\mathbf N_y)\in I_y)=1-O(L^{-1/5}),
\]
uniformly for the two indicated values of $y$.
\end{proposition}
\begin{proof}
Since $\mathbf N_y\leq x^{\alpha^3}$, \eqref{eq:orbit-log} gives
\[
 \log\frac{\Syr^{n_0}(\mathbf N_y)}x
 \leq(\alpha^3-1)L-\beta n_0+C_0L^{3/5}<0
\]
for large $x$.  For example $\alpha^3-1<1/300$ and
$\beta/(10\log2)>1/40$, whereas the floor costs only a constant.
Since $\mathbf N_y>x$, the first crossing is therefore at an integer
$1\leq t\leq n_0$.  Applying \eqref{eq:orbit-log} at $t$ and $t-1$
gives
\begin{equation}\label{eq:time-log}
 \left|t-\frac{\log(\mathbf N_y/x)}\beta\right|
 \leq C_1L^{3/5}.
\end{equation}
Indeed the inequalities $\Syr^t(\mathbf N_y)\leq x$
and $\Syr^{t-1}(\mathbf N_y)>x$ bound $t$ from below and above;
the additional one-step term is absorbed into $C_1$.

If $\mathbf N_y\in J_y$, the central expression in
\eqref{eq:time-log} lies inside $I_y$ with extra margin
$(\beta^{-1}-1)L^{4/5}$ on both sides.  Since $0<\beta<1$, this
dominates $C_1L^{3/5}$.  This proves the event inclusion.

For $1\leq a<b$, integral comparison along the odd progression gives
\[
 \sum_{\substack{a\leq m\leq b\\m\text{ odd}}}\frac1m
 =\frac12\log(b/a)+O(1/a),
\]
uniformly in the real endpoints.  Thus the block denominator is
$\frac12(\alpha-1)\log y+O(1/y)$, whereas the two excluded pieces have
total harmonic mass $L^{4/5}+O(1/y)$.  Their probability is
$O(L^{-1/5})$.  The logarithmic width of $J_y$ is
$(\alpha-1)\log y-2L^{4/5}>0$.  Finally add
$\Prob(G_x^c)=O(L^{-10})$.
\end{proof}

Thus the multiplicative buffer gives the same interval $I_y$ and
exceptional probability $O(L^{-1/5})$ required in the subsequent argument.

\subsection{A uniform coefficient bound over valuation words}

For a positive tuple $\mathbf a=(a_1,\ldots,a_m)$ put
\[
 A=a_1+\cdots+a_m,\qquad b=a_m,\qquad
 F_m(\mathbf a)=\sum_{i=1}^m3^{m-i}2^{-(a_i+\cdots+a_m)},
 \quad F_0=0.
\]
Composition of the maps $z\mapsto(3z+1)/2^{a_i}$ gives
\begin{equation}\label{eq:affine}
 \Syr^m(M)=3^m2^{-A}M+F_m(\mathbf a)
\end{equation}
when $\mathbf a$ is the actual valuation tuple of $M$.
Induction proves the formula: the next map multiplies every existing term
by $3/2^{a_{m+1}}$ and adds $2^{-a_{m+1}}$.
In particular $0\leq F_m\leq3^m$, a bound uniform in the tuple.
For the latter bound, replace every negative power of two by at most
$1/2$ and sum the finite geometric series.

For an arbitrary $E\subseteq\Opos\cap[1,x]$ set
\[
 E'=\left\{M\in\Opos:\begin{array}{l}
 t_x(M)=m_0,\quad p_x(M)\in E,\\
 e^{-L^{7/10}}(4/3)^{m_0}x\leq M
       \leq e^{L^{7/10}}(4/3)^{m_0}x
 \end{array}\right\}.
\]
This is the same set as in the approximate formula of
\cite[Proposition 5.2]{Tao7}; successful passage makes its meaning
independent of the failure convention.  For $n\in I_y\cap\Z$ write
$m=m_0$, $r=n-m\geq0$, and define on $\Z/3^r\Z$
\begin{equation}\label{eq:cn}
 c_n(z)=3^r\sum_{\substack{M\in E'\\M\equiv z\pmod{3^r}}}\frac1M.
\end{equation}
The source parameters ensure $m\leq n\leq n_0$ for sufficiently large
$x$.  The case $r=0$ is allowed in the estimate below, with modulus one.

\begin{proposition}[Uniform coefficient bound]\label{prop:cn}
There is an absolute constant $C$ such that $0\leq c_n(z)\leq C$
for all sufficiently large $x$, both choices of $y$, all
$n\in I_y\cap\Z$, all residue classes $z$, and every set $E$ above.
\end{proposition}
\begin{proof}
Partition \eqref{eq:cn} by the unique length-$m$ valuation tuple of $M$,
and call a tuple's contribution $c_{n,\mathbf a}(z)$.  At its first
crossing, the exact affine formulas give
\[
 3^m2^{-A}M+F_m\leq x
 <3^{m-1}2^{-(A-b)}M+F_{m-1}.
\]
The coefficient $3^{m-1}$ corresponds to the preceding $m-1$ iterates.
Uniformly $F_{m-1}\leq3^{m-1}=o(x)$, so every such $M$ lies between
\begin{equation}\label{eq:M-bounds}
 M_0=3^{-m}2^{A-b}x\quad\hbox{and}\quad M_1=3^{-m}2^A x
\end{equation}
for large $x$.  Indeed the upper bound follows by discarding $F_m\geq0$,
and the lower follows from $x-F_{m-1}>x/3$.

The integer $B_{\mathbf a}=2^A F_m$ is odd: its first term is
$3^{m-1}$ and all its later terms are even.  The terminal value in
\eqref{eq:affine} is odd, hence
\[
 3^mM+B_{\mathbf a}\equiv2^A\pmod{2^{A+1}}.
\]
Since $3^m$ is a unit modulo $2^{A+1}$, this requires a single class of
$M$ modulo $2^{A+1}$.  Combining it with $M\equiv z\pmod{3^r}$ gives
a single class modulo $q=2^{A+1}3^r$ by the Chinese remainder theorem.
Only this necessary restriction is needed for the upper bound.

For any progression and $1\leq U\leq V$,
\begin{equation}\label{eq:progression}
 \sum_{\substack{U\leq M\leq V\\M\equiv z\pmod q}}\frac1M
 \leq\frac1U+\frac1q\log(V/U).
\end{equation}
To prove it, keep the first progression term separately; for each
subsequent term at $t$, use $1/t\leq q^{-1}\int_{t-q}^{t}du/u$.
The integration intervals are disjoint and lie inside $[U,V]$.
If there are no progression terms, the bound is immediate.

If $2^A\leq x^{1/2}$, then
\[
 \frac q{M_0}=\frac{2^{b+1}3^n}{x}\leq2x^{-17/50}\leq1
\]
for large $x$.  Here $b\leq A$ and $3^n\leq x^{4/25}$, since
$n\leq L/(10\log2)$ and $3^5<2^8$.
Because $M_1/M_0=2^b$ and $b\geq1$, \eqref{eq:progression} gives
\begin{equation}\label{eq:low}
 c_{n,\mathbf a}(z)\leq C_2 b2^{-A}.
\end{equation}

If $2^A>x^{1/2}$, set $V_*=3^{m-1}2^{-(A-b)}M$.
The previous crossing inequality gives $V_*>x-F_{m-1}$, and therefore
$F_{m-1}=o(V_*)$ uniformly.  The last valuation is
\[
 b=\vnu(3(V_*+F_{m-1})+1).
\]
The positive integer inside the valuation is divisible by $2^b$, so
$2^b\leq3(V_*+F_{m-1})+1\leq5V_*$ for large $x$.
Consequently $M\geq(3/5)3^{-m}2^A$.  Apply
\eqref{eq:progression} with lower endpoint
$\max(M_0,(3/5)3^{-m}2^A)$ and upper endpoint $M_1$.
If they are out of order the contribution is zero; otherwise their ratio
is at most $2^b$.  We obtain
\begin{equation}\label{eq:high}
 c_{n,\mathbf a}(z)\leq C_3(3^n2^{-A}+b2^{-A}).
\end{equation}

The weights must now be summed before weakening them.  Exactly,
\[
 \sum_{a_1,\ldots,a_m\geq1}b2^{-A}
 =\left(\sum_{a\geq1}2^{-a}\right)^{m-1}
    \sum_{b\geq1}b2^{-b}=2.
\]
For the additional term in \eqref{eq:high},
\begin{align*}
 3^n\sum_{2^A>x^{1/2}}2^{-A}
 &\leq3^nx^{-1/4}\sum_{a_1,\ldots,a_m\geq1}2^{-A/2}\\
 &=3^nx^{-1/4}(1+\sqrt2)^m
 \leq x^{4/25-1/4+1/100000}
 =x^{-8999/100000}.
\end{align*}
The last step uses $1+\sqrt2<e$ and $m\leq L/100000$.
All summands are nonnegative, so enlarging the actual finite tuple range
to all positive tuples is valid.  Summing \eqref{eq:low} and
\eqref{eq:high} proves $c_n(z)\leq C_4(2+x^{-8999/100000})$, uniformly
in every stated variable.
\end{proof}

The last displayed majorant in the printed proof of Lemma~5.3 is instead
$\sum_{a_1,\ldots,a_m\geq1}2^{-A/2}=(1+\sqrt2)^m$.
It is finite for a fixed $m$ but is not uniformly bounded when
$m=\lfloor L/100000\rfloor$ grows.  Keeping the sharper weights already
available in the preceding estimates proves exactly the uniform result
needed by that lemma.

To apply the uniform bound, first observe that the source range gives
$r=n-m\geq m$: the lower endpoint gives
$n\geq(\alpha-1)L/\beta+L^{4/5}$, whereas
$2m\leq(\alpha-1)L/50$ and $\beta<1$.
If $p_r$ is the law of Tao's
$r$-Syracuse random variable on $\Z/3^r\Z$, define the uniform lift of
its reduction to $\Z/3^m\Z$ by
\[
 (L_{m\to r}p_m)(z)=3^{m-r}p_m(z\bmod3^m)\quad(0\leq m\leq r).
\]
Reduction has fibres of size $3^{r-m}$, so this lift preserves mass.
Tao's projective identity says that reducing $p_r$ gives $p_m$.
His fine-scale mixing estimate controls
$\normone{p_r-L_{m\to r}p_m}$.  Proposition~\ref{prop:cn} now gives
the precise dual estimate
\[
 \left|\sum_{z\bmod3^r}c_n(z)
       \bigl(p_r(z)-(L_{m\to r}p_m)(z)\bigr)\right|
 \leq C\normone{p_r-L_{m\to r}p_m}.
\]
The coefficient therefore transfers the fine-scale mixing error with
a constant independent of $x$, $n$, and $E$.
